% Options for packages loaded elsewhere
\PassOptionsToPackage{unicode}{hyperref}
\PassOptionsToPackage{hyphens}{url}
\documentclass[
]{article}
\usepackage{xcolor}
\usepackage[margin=1in]{geometry}
\usepackage{amsmath,amssymb}
\setcounter{secnumdepth}{-\maxdimen} % remove section numbering
\usepackage{iftex}
\ifPDFTeX
  \usepackage[T1]{fontenc}
  \usepackage[utf8]{inputenc}
  \usepackage{textcomp} % provide euro and other symbols
\else % if luatex or xetex
  \usepackage{unicode-math} % this also loads fontspec
  \defaultfontfeatures{Scale=MatchLowercase}
  \defaultfontfeatures[\rmfamily]{Ligatures=TeX,Scale=1}
\fi
\usepackage{lmodern}
\ifPDFTeX\else
  % xetex/luatex font selection
\fi
% Use upquote if available, for straight quotes in verbatim environments
\IfFileExists{upquote.sty}{\usepackage{upquote}}{}
\IfFileExists{microtype.sty}{% use microtype if available
  \usepackage[]{microtype}
  \UseMicrotypeSet[protrusion]{basicmath} % disable protrusion for tt fonts
}{}
\makeatletter
\@ifundefined{KOMAClassName}{% if non-KOMA class
  \IfFileExists{parskip.sty}{%
    \usepackage{parskip}
  }{% else
    \setlength{\parindent}{0pt}
    \setlength{\parskip}{6pt plus 2pt minus 1pt}}
}{% if KOMA class
  \KOMAoptions{parskip=half}}
\makeatother
\usepackage{color}
\usepackage{fancyvrb}
\newcommand{\VerbBar}{|}
\newcommand{\VERB}{\Verb[commandchars=\\\{\}]}
\DefineVerbatimEnvironment{Highlighting}{Verbatim}{commandchars=\\\{\}}
% Add ',fontsize=\small' for more characters per line
\usepackage{framed}
\definecolor{shadecolor}{RGB}{248,248,248}
\newenvironment{Shaded}{\begin{snugshade}}{\end{snugshade}}
\newcommand{\AlertTok}[1]{\textcolor[rgb]{0.94,0.16,0.16}{#1}}
\newcommand{\AnnotationTok}[1]{\textcolor[rgb]{0.56,0.35,0.01}{\textbf{\textit{#1}}}}
\newcommand{\AttributeTok}[1]{\textcolor[rgb]{0.13,0.29,0.53}{#1}}
\newcommand{\BaseNTok}[1]{\textcolor[rgb]{0.00,0.00,0.81}{#1}}
\newcommand{\BuiltInTok}[1]{#1}
\newcommand{\CharTok}[1]{\textcolor[rgb]{0.31,0.60,0.02}{#1}}
\newcommand{\CommentTok}[1]{\textcolor[rgb]{0.56,0.35,0.01}{\textit{#1}}}
\newcommand{\CommentVarTok}[1]{\textcolor[rgb]{0.56,0.35,0.01}{\textbf{\textit{#1}}}}
\newcommand{\ConstantTok}[1]{\textcolor[rgb]{0.56,0.35,0.01}{#1}}
\newcommand{\ControlFlowTok}[1]{\textcolor[rgb]{0.13,0.29,0.53}{\textbf{#1}}}
\newcommand{\DataTypeTok}[1]{\textcolor[rgb]{0.13,0.29,0.53}{#1}}
\newcommand{\DecValTok}[1]{\textcolor[rgb]{0.00,0.00,0.81}{#1}}
\newcommand{\DocumentationTok}[1]{\textcolor[rgb]{0.56,0.35,0.01}{\textbf{\textit{#1}}}}
\newcommand{\ErrorTok}[1]{\textcolor[rgb]{0.64,0.00,0.00}{\textbf{#1}}}
\newcommand{\ExtensionTok}[1]{#1}
\newcommand{\FloatTok}[1]{\textcolor[rgb]{0.00,0.00,0.81}{#1}}
\newcommand{\FunctionTok}[1]{\textcolor[rgb]{0.13,0.29,0.53}{\textbf{#1}}}
\newcommand{\ImportTok}[1]{#1}
\newcommand{\InformationTok}[1]{\textcolor[rgb]{0.56,0.35,0.01}{\textbf{\textit{#1}}}}
\newcommand{\KeywordTok}[1]{\textcolor[rgb]{0.13,0.29,0.53}{\textbf{#1}}}
\newcommand{\NormalTok}[1]{#1}
\newcommand{\OperatorTok}[1]{\textcolor[rgb]{0.81,0.36,0.00}{\textbf{#1}}}
\newcommand{\OtherTok}[1]{\textcolor[rgb]{0.56,0.35,0.01}{#1}}
\newcommand{\PreprocessorTok}[1]{\textcolor[rgb]{0.56,0.35,0.01}{\textit{#1}}}
\newcommand{\RegionMarkerTok}[1]{#1}
\newcommand{\SpecialCharTok}[1]{\textcolor[rgb]{0.81,0.36,0.00}{\textbf{#1}}}
\newcommand{\SpecialStringTok}[1]{\textcolor[rgb]{0.31,0.60,0.02}{#1}}
\newcommand{\StringTok}[1]{\textcolor[rgb]{0.31,0.60,0.02}{#1}}
\newcommand{\VariableTok}[1]{\textcolor[rgb]{0.00,0.00,0.00}{#1}}
\newcommand{\VerbatimStringTok}[1]{\textcolor[rgb]{0.31,0.60,0.02}{#1}}
\newcommand{\WarningTok}[1]{\textcolor[rgb]{0.56,0.35,0.01}{\textbf{\textit{#1}}}}
\usepackage{longtable,booktabs,array}
\newcounter{none} % for unnumbered tables
\usepackage{calc} % for calculating minipage widths
% Correct order of tables after \paragraph or \subparagraph
\usepackage{etoolbox}
\makeatletter
\patchcmd\longtable{\par}{\if@noskipsec\mbox{}\fi\par}{}{}
\makeatother
% Allow footnotes in longtable head/foot
\IfFileExists{footnotehyper.sty}{\usepackage{footnotehyper}}{\usepackage{footnote}}
\makesavenoteenv{longtable}
\usepackage{graphicx}
\makeatletter
\newsavebox\pandoc@box
\newcommand*\pandocbounded[1]{% scales image to fit in text height/width
  \sbox\pandoc@box{#1}%
  \Gscale@div\@tempa{\textheight}{\dimexpr\ht\pandoc@box+\dp\pandoc@box\relax}%
  \Gscale@div\@tempb{\linewidth}{\wd\pandoc@box}%
  \ifdim\@tempb\p@<\@tempa\p@\let\@tempa\@tempb\fi% select the smaller of both
  \ifdim\@tempa\p@<\p@\scalebox{\@tempa}{\usebox\pandoc@box}%
  \else\usebox{\pandoc@box}%
  \fi%
}
% Set default figure placement to htbp
\def\fps@figure{htbp}
\makeatother
% definitions for citeproc citations
\NewDocumentCommand\citeproctext{}{}
\NewDocumentCommand\citeproc{mm}{%
  \begingroup\def\citeproctext{#2}\cite{#1}\endgroup}
\makeatletter
 % allow citations to break across lines
 \let\@cite@ofmt\@firstofone
 % avoid brackets around text for \cite:
 \def\@biblabel#1{}
 \def\@cite#1#2{{#1\if@tempswa , #2\fi}}
\makeatother
\newlength{\cslhangindent}
\setlength{\cslhangindent}{1.5em}
\newlength{\csllabelwidth}
\setlength{\csllabelwidth}{3em}
\newenvironment{CSLReferences}[2] % #1 hanging-indent, #2 entry-spacing
 {\begin{list}{}{%
  \setlength{\itemindent}{0pt}
  \setlength{\leftmargin}{0pt}
  \setlength{\parsep}{0pt}
  % turn on hanging indent if param 1 is 1
  \ifodd #1
   \setlength{\leftmargin}{\cslhangindent}
   \setlength{\itemindent}{-1\cslhangindent}
  \fi
  % set entry spacing
  \setlength{\itemsep}{#2\baselineskip}}}
 {\end{list}}
\usepackage{calc}
\newcommand{\CSLBlock}[1]{\hfill\break\parbox[t]{\linewidth}{\strut\ignorespaces#1\strut}}
\newcommand{\CSLLeftMargin}[1]{\parbox[t]{\csllabelwidth}{\strut#1\strut}}
\newcommand{\CSLRightInline}[1]{\parbox[t]{\linewidth - \csllabelwidth}{\strut#1\strut}}
\newcommand{\CSLIndent}[1]{\hspace{\cslhangindent}#1}
\setlength{\emergencystretch}{3em} % prevent overfull lines
\providecommand{\tightlist}{%
  \setlength{\itemsep}{0pt}\setlength{\parskip}{0pt}}
\usepackage{hyperref}
\usepackage{bookmark}
\IfFileExists{xurl.sty}{\usepackage{xurl}}{} % add URL line breaks if available
\urlstyle{same}
\hypersetup{
  hidelinks,
  pdfcreator={LaTeX via pandoc}}

\author{}
\date{\vspace{-2.5em}}

\begin{document}

\begin{center}
{\LARGE \textbf{gp3tools: Reproducible Analysis and Reporting of Gazepoint GP3 Eye-Tracking Exports}\\}
\vspace{1em}
Stefanos Balaskas\\
Department of Business Administration, University of Patras, Patras, Greece\\
Correspondence: \href{mailto:s.balaskas@ac.upatras.gr}{s.balaskas@ac.upatras.gr}\\
ORCID: \href{https://orcid.org/0000-0003-2444-9796}{https://orcid.org/0000-0003-2444-9796}
\end{center}

\begin{abstract}
Eye-tracking studies increasingly require reproducible software pipelines that connect device exports with transparent preprocessing, quality control, statistical analysis, visualization, and reporting. Gazepoint GP3 and Gazepoint Analysis exports provide accessible gaze, fixation, pupil, timing, media, and area-of-interest data, but exported CSV folders usually require systematic validation before they can support defensible inference. \texttt{gp3tools} is an R package for importing, validating, preprocessing, modelling, visualising, and reporting Gazepoint GP3 / Gazepoint Analysis exports. The package supports file discovery, file-pair validation, sampling diagnostics, tracking-quality summaries, pupil preprocessing, area-of-interest summaries, fixation and transition workflows, model-ready data preparation, diagnostic plotting, exclusion recommendations, reporting checklists, and reproducible report generation. This paper presents the motivation, design, validation strategy, and synthetic demonstration workflow of \texttt{gp3tools}. The demonstration is used to evaluate workflow reproducibility and software behaviour; it is not used to draw substantive psychological, perceptual, or marketing conclusions. The package fills a practical gap between vendor-specific Gazepoint CSV exports and publication-ready R workflows for transparent eye-tracking analysis.
\end{abstract}

\section{Introduction}\label{introduction}

Eye-tracking research increasingly depends on reproducible software
pipelines that connect raw device exports with transparent
preprocessing, statistical analysis, visualization, and reporting. These
pipelines matter because gaze, fixation, pupil, and area-of-interest
(AOI) outcomes are not raw facts extracted directly from an eye tracker.
They are constructed through device export formats, calibration quality,
validity flags, timing fields, preprocessing rules, AOI definitions,
trial windows, model specifications, and reporting decisions.
Reproducibility therefore depends not only on the final statistical
model but also on the traceability of the steps that precede it
(\citeproc{ref-Wilson2014}{Wilson et al. 2014}).

Gazepoint GP3 and Gazepoint Analysis are used in experimental and
applied eye-tracking contexts because they provide accessible export
workflows for gaze, fixation, pupil, timing, media, and AOI information
(\citeproc{ref-GazepointGP3}{Gazepoint Research Inc. 2026a},
\citeproc{ref-GazepointAPI}{2026b}). However, exported CSV files are not
automatically equivalent to analysis-ready data. A typical project may
require checking whether all-gaze and fixation files are correctly
paired, whether sampling rate is plausible, whether gaze and pupil
validity are adequate, whether AOI labels and denominators are
consistent, whether pupil artifacts and baseline decisions have been
handled transparently, and whether the final outputs are sufficiently
documented for reporting.

These requirements are not merely technical. Preprocessing choices can
affect downstream interpretation. Pupil data are sensitive to blinks,
artifacts, interpolation, and baseline correction
(\citeproc{ref-KretSjakShie2019}{Kret and Sjak-Shie 2019};
\citeproc{ref-Mathot2018Baseline}{Mathot et al. 2018};
\citeproc{ref-Mathot2018Pupillometry}{Mathot 2018}). AOI analyses depend
on how areas of interest are defined, interpreted, and reported
(\citeproc{ref-Orquin2016}{Orquin et al. 2016}). More generally,
research-software practices emphasize version control, testing,
documentation, provenance, and citable software releases as part of
reproducible computational work (\citeproc{ref-Wilson2014}{Wilson et al.
2014}; \citeproc{ref-Smith2016}{Smith et al. 2016};
\citeproc{ref-Katz2021}{Katz et al. 2021}).

\texttt{gp3tools} was developed to make Gazepoint GP3 / Gazepoint
Analysis workflows more explicit and reproducible in R
(\citeproc{ref-RCoreTeam2026}{R Core Team 2026}). The package does not
propose a new statistical estimator, eye-tracking theory, or universal
preprocessing standard. Its contribution is a Gazepoint-specific
workflow layer that converts vendor-style export folders into auditable
R objects, diagnostic summaries, model-ready tables, plots, reports, and
reporting-readiness outputs.

\section{Statement of need}\label{statement-of-need}

Several tools already support eye-tracking and pupillometry workflows.
\texttt{eyetrackingR} provides R-based tools for preparing, summarizing,
visualizing, and analysing eye-tracking data
(\citeproc{ref-DinkFerguson2015}{Dink and Ferguson 2015}).
\texttt{PupillometryR} provides a pipeline for cleaning, preparing,
plotting, and analysing pupillometry data
(\citeproc{ref-Forbes2020}{Forbes 2020}). General R packages support
reproducible documents, statistical graphics, mixed models, generalized
additive models, and generalized linear mixed models
(\citeproc{ref-Wickham2016ggplot2}{Wickham 2016};
\citeproc{ref-Xie2015knitr}{Xie 2015};
\citeproc{ref-XieAllaireGrolemund2018}{Xie et al. 2018};
\citeproc{ref-Bates2015}{Bates et al. 2015};
\citeproc{ref-Brooks2017}{Brooks et al. 2017};
\citeproc{ref-Wood2017}{Wood 2017}).

The need for \texttt{gp3tools} is therefore specific rather than
generic. Existing tools are often platform-agnostic, signal-specific,
paradigm-specific, or focused on later analysis stages. Gazepoint users
need an export-aware workflow that starts from the concrete structure of
Gazepoint GP3 / Gazepoint Analysis CSV folders and proceeds through
validation, preprocessing, modelling, visualization, and reporting in
the same R environment.

\texttt{gp3tools} addresses this gap by providing a reusable,
documented, and testable workflow for Gazepoint exports. It is intended
for researchers who need to process private Gazepoint recordings while
sharing transparent code, synthetic demonstration data, diagnostic
summaries, and reporting-ready outputs. Its intended users include
applied eye-tracking researchers, behavioural scientists, marketing and
consumer researchers, human-computer interaction researchers, usability
laboratories, teaching laboratories, and methodologically oriented users
who need auditable preprocessing before statistical modelling.

\section{Software contribution}\label{software-contribution}

\texttt{gp3tools} makes three software contributions. First, it provides
a Gazepoint-specific import and validation layer for participant-level
all-gaze, fixation, and summary exports. Second, it provides
audit-oriented workflows for sampling, tracking quality, pupil
preprocessing, AOI summaries, transition structures, modelling
preparation, reporting readiness, and exclusion recommendations. Third,
it provides public reproducibility material, including synthetic
Gazepoint-style exports, R Markdown workflows, vignettes, tests, and a
Zenodo-archived software release
(\citeproc{ref-Balaskas2026gp3tools}{Balaskas 2026}).

The package is best understood as workflow infrastructure. It supports
the data-engineering, validation, and reporting stages that precede
statistical inference. This distinction is important because many
eye-tracking results depend on preprocessing and quality-control
decisions that are often described only briefly in empirical
manuscripts. By making these decisions explicit and reproducible,
\texttt{gp3tools} helps researchers document how raw-style Gazepoint
exports become analysis-ready datasets.

\section{Package design}\label{package-design}

\texttt{gp3tools} is organized around workflow families rather than
isolated utilities. The first layer handles file discovery, import, and
harmonization. It reads Gazepoint export folders, identifies all-gaze,
fixation, and summary files, checks file pairing, and creates structured
sample-level and fixation-level objects. This layer reduces the risk
that later analyses are based on incomplete, duplicate, or mismatched
files.

The second layer handles diagnostics and quality control. It estimates
sampling rate, summarizes tracking quality, evaluates gaze and pupil
validity, flags recordings requiring review, checks master-table
structure, audits AOI coding, evaluates condition-level quality
imbalance, and supports gaze recalibration when known target coordinates
are available. These outputs are intended to be inspectable artifacts
rather than hidden preprocessing steps.

The third layer prepares analysis-ready summaries. It supports
AOI-window summaries, pupil-window summaries, fixation-level summaries,
fixation-aligned data, event-aligned workflows, AOI transition tables,
transition matrices, Markov-style objects, hidden-Markov-model-ready
tables, semi-Markov-ready tables, time-course data, growth-curve-style
data, and confirmatory model-ready objects for mixed-effects modelling.

The fourth layer supports modelling, robustness, reporting, and
interoperability. It includes preparation functions or wrappers for
linear mixed models, generalized linear mixed models, GAMM-style
workflows, nested-model comparison, prediction plotting,
residual-diagnostic workflows, leave-one-unit sensitivity,
divergence-point estimation, multiverse-style preprocessing summaries,
reporting checklists, exclusion recommendations, real-data readiness
gates, manuscript-ready model tables, and adapters for related
eye-tracking packages.

\subsection{Workflow map}\label{workflow-map}

Table 1 provides a compact map of the main workflow families. The full
exported-function inventory remains in Supplementary Material S1.

\begin{longtable}[]{@{}
  >{\raggedright\arraybackslash}p{(\linewidth - 4\tabcolsep) * \real{0.1889}}
  >{\raggedright\arraybackslash}p{(\linewidth - 4\tabcolsep) * \real{0.3833}}
  >{\raggedright\arraybackslash}p{(\linewidth - 4\tabcolsep) * \real{0.4278}}@{}}
\caption{Main workflow families supported by gp3tools.}\tabularnewline
\toprule\noalign{}
\begin{minipage}[b]{\linewidth}\raggedright
Workflow family
\end{minipage} & \begin{minipage}[b]{\linewidth}\raggedright
Purpose
\end{minipage} & \begin{minipage}[b]{\linewidth}\raggedright
Example outputs
\end{minipage} \\
\midrule\noalign{}
\endfirsthead
\toprule\noalign{}
\begin{minipage}[b]{\linewidth}\raggedright
Workflow family
\end{minipage} & \begin{minipage}[b]{\linewidth}\raggedright
Purpose
\end{minipage} & \begin{minipage}[b]{\linewidth}\raggedright
Example outputs
\end{minipage} \\
\midrule\noalign{}
\endhead
\bottomrule\noalign{}
\endlastfoot
Import and file inspection & Read Gazepoint-style exports and identify
expected file types. & classified files; imported all-gaze, fixation,
and summary tables \\
Master-table validation & Create and inspect analysis-ready combined
data. & validated master tables; column diagnostics; duplicate-key
checks \\
Sampling and tracking quality & Check sampling behaviour, gaze validity,
and review flags. & sampling-rate summaries; tracking-quality tables;
review flags \\
Pupil preprocessing & Prepare pupil data for window and time-course
analysis. & baseline-corrected pupil series; interpolation flags;
pupil-window summaries \\
AOI workflows & Summarise and audit gaze behaviour by area of interest.
& AOI sample and fixation summaries; transition matrices; denominator
audits \\
Model preparation and sensitivity & Prepare model-ready datasets and
evaluate robustness. & GLMM/GAMM input tables; leave-one-out summaries;
divergence-point results \\
Reporting and export & Generate reporting checklists, exports, and
reproducibility outputs. & reporting checklists; exclusion
recommendations; exported CSV tables \\
\end{longtable}

\section{Demonstration case study}\label{demonstration-case-study}

The demonstration uses a synthetic Gazepoint-style export dataset. The
purpose is to exercise the package under realistic but shareable
conditions, not to test a substantive psychological or marketing
hypothesis. The synthetic data include all-gaze files, fixation files, a
summary export, media identifiers, AOI states, gaze coordinates, pupil
values, validity indicators, missing pupil intervals, invalid samples,
and participant-level quality variation.

The case study demonstrates that \texttt{gp3tools} can identify complete
all-gaze/fixation file pairs, estimate sampling rate, flag
review-required recordings, create analysis-ready all-gaze and fixation
tables, summarize AOI viewing behaviour, audit AOI denominators,
preprocess pupil data, run interpolation and baseline-quality checks,
prepare confirmatory pupil-window and AOI-window model data, construct
transition and sequence summaries, create exclusion recommendations,
generate reporting-readiness objects, and produce diagnostic report
outputs.

\subsection{Minimal reproducible
workflow}\label{minimal-reproducible-workflow}

The full executable workflow is provided in Supplementary Material S1.
The main manuscript reports the minimal code structure used by a typical
user.

\begin{Shaded}
\begin{Highlighting}[]
\CommentTok{\# Development version}
\NormalTok{remotes}\SpecialCharTok{::}\FunctionTok{install\_github}\NormalTok{(}\StringTok{"stefanosbalaskas/gp3tools"}\NormalTok{)}

\CommentTok{\# After CRAN acceptance}
\CommentTok{\# install.packages("gp3tools")}
\end{Highlighting}
\end{Shaded}

\begin{Shaded}
\begin{Highlighting}[]
\FunctionTok{library}\NormalTok{(gp3tools)}

\NormalTok{demo\_dir }\OtherTok{\textless{}{-}} \FunctionTok{system.file}\NormalTok{(}
  \StringTok{"extdata"}\NormalTok{,}
  \StringTok{"gazepoint\_realistic\_demo\_exports"}\NormalTok{,}
  \AttributeTok{package =} \StringTok{"gp3tools"}
\NormalTok{)}

\NormalTok{workflow }\OtherTok{\textless{}{-}} \FunctionTok{run\_gazepoint\_workflow}\NormalTok{(}
  \AttributeTok{input\_dir =}\NormalTok{ demo\_dir,}
  \AttributeTok{output\_dir =} \StringTok{"gp3tools\_demo\_output"}\NormalTok{,}
  \AttributeTok{prefix =} \StringTok{"gp3tools\_demo"}\NormalTok{,}
  \AttributeTok{save\_plots =} \ConstantTok{TRUE}\NormalTok{,}
  \AttributeTok{create\_report =} \ConstantTok{TRUE}
\NormalTok{)}

\FunctionTok{summarise\_gazepoint\_workflow}\NormalTok{(workflow)}
\end{Highlighting}
\end{Shaded}

\subsection{Demonstration outputs}\label{demonstration-outputs}

The larger paper-only synthetic showcase contains 48 synthetic all-gaze
files, 48 synthetic fixation files, and one summary file. The end-to-end
run produces 115,200 all-gaze rows, 5,691 fixation rows, 192
sampling-quality rows, 1,344 AOI summary rows, and 60 review-required
rows. All 48 expected file pairs are complete. These values are
software-demonstration outputs; they are not interpreted as empirical
findings.

\begin{longtable}[]{@{}
  >{\raggedright\arraybackslash}p{(\linewidth - 4\tabcolsep) * \real{0.2619}}
  >{\raggedleft\arraybackslash}p{(\linewidth - 4\tabcolsep) * \real{0.0833}}
  >{\raggedright\arraybackslash}p{(\linewidth - 4\tabcolsep) * \real{0.6548}}@{}}
\caption{Summary of synthetic demonstration outputs.}\tabularnewline
\toprule\noalign{}
\begin{minipage}[b]{\linewidth}\raggedright
Output
\end{minipage} & \begin{minipage}[b]{\linewidth}\raggedleft
Count
\end{minipage} & \begin{minipage}[b]{\linewidth}\raggedright
Interpretation
\end{minipage} \\
\midrule\noalign{}
\endfirsthead
\toprule\noalign{}
\begin{minipage}[b]{\linewidth}\raggedright
Output
\end{minipage} & \begin{minipage}[b]{\linewidth}\raggedleft
Count
\end{minipage} & \begin{minipage}[b]{\linewidth}\raggedright
Interpretation
\end{minipage} \\
\midrule\noalign{}
\endhead
\bottomrule\noalign{}
\endlastfoot
All-gaze rows & 115200 & sample-level gaze and pupil observations \\
Fixation rows & 5691 & fixation-level observations \\
Sampling-quality rows & 192 & participant-by-media sampling
diagnostics \\
AOI summary rows & 1344 & area-of-interest summaries across
trials/windows \\
Review-required rows & 60 & rows flagged for review rather than
automatic deletion \\
Complete file pairs & 48 & all expected all-gaze/fixation pairs found \\
\end{longtable}

\begin{figure}
\centering
\pandocbounded{\includegraphics[keepaspectratio,alt={Scale of the main synthetic demonstration outputs. Counts are shown on a log10 scale because the outputs differ by several orders of magnitude.}]{C:/Users/Stefanos-PC/Documents/Rstudio/gp3tools/paper-showcase/arxiv/gp3tools_arxiv_files/ab6d35b591c2d35ee166f65bdd72ad9506a53be3.pdf}}
\caption{Scale of the main synthetic demonstration outputs. Counts are
shown on a log10 scale because the outputs differ by several orders of
magnitude.}
\end{figure}

\section{Validation and software
quality}\label{validation-and-software-quality}

The validation strategy combines automated tests, reproducible examples,
synthetic demonstration data, package checks, documentation, and release
archiving. This strategy follows general principles of scientific
computing, including readable code, version control, automated testing,
scripted workflows, and preserved provenance
(\citeproc{ref-Wilson2014}{Wilson et al. 2014}). It also aligns with
software-citation and FAIR-oriented practices through public metadata,
versioned releases, persistent citation information, and reusable
documentation (\citeproc{ref-Smith2016}{Smith et al. 2016};
\citeproc{ref-Katz2021}{Katz et al. 2021}).

The package includes unit tests for core functions and workflow
branches, following the standard R testing approach supported by
\texttt{testthat} (\citeproc{ref-Wickham2011testthat}{Wickham 2011}).
The test suite checks import behaviour, validation outputs, diagnostic
summaries, pupil preprocessing, AOI workflows, transition workflows,
modelling preparation, reporting utilities, optional-dependency
handling, and structured error/status reporting.

\begin{longtable}[]{@{}
  >{\raggedright\arraybackslash}p{(\linewidth - 2\tabcolsep) * \real{0.2017}}
  >{\raggedright\arraybackslash}p{(\linewidth - 2\tabcolsep) * \real{0.7983}}@{}}
\caption{Software-quality evidence reported for
gp3tools.}\tabularnewline
\toprule\noalign{}
\begin{minipage}[b]{\linewidth}\raggedright
Evidence
\end{minipage} & \begin{minipage}[b]{\linewidth}\raggedright
Role
\end{minipage} \\
\midrule\noalign{}
\endfirsthead
\toprule\noalign{}
\begin{minipage}[b]{\linewidth}\raggedright
Evidence
\end{minipage} & \begin{minipage}[b]{\linewidth}\raggedright
Role
\end{minipage} \\
\midrule\noalign{}
\endhead
\bottomrule\noalign{}
\endlastfoot
Automated tests & Checks core functions, edge cases, optional-dependency
paths, and workflow objects. \\
Package checks & Confirms package installation, documentation, examples,
vignettes, and R CMD check compliance. \\
Cross-platform CI & Evaluates package behaviour across major operating
systems and R-devel/release settings. \\
Release archiving & Provides a persistent versioned software citation
through Zenodo. \\
Documentation & Provides README guidance, function help, vignettes,
example data, and case-study workflows. \\
Synthetic demonstration & Exercises the export-to-report workflow
without exposing private participant data. \\
\end{longtable}

The software-quality claims in this paper are intentionally narrow. The
demonstration shows that the package can execute a reproducible workflow
on Gazepoint-like exports and produce expected diagnostic, analysis, and
reporting objects. It does not prove that all possible Gazepoint exports
will require no local inspection, nor does it validate any psychological
theory. It also does not independently validate vendor hardware
specifications such as nominal sampling rate or accuracy; those remain
vendor specifications rather than package-level validation results.

\section{Data and code availability}\label{data-and-code-availability}

The \texttt{gp3tools} source code is publicly available at
\texttt{https://github.com/stefanosbalaskas/gp3tools}. Version 1.0.1 has
been archived on Zenodo and should be cited using DOI
\texttt{10.5281/zenodo.20731340}
(\citeproc{ref-Balaskas2026gp3tools}{Balaskas 2026}). At the time of
writing, the package has been submitted to CRAN and is awaiting review.

The repository includes source code, documentation, tests, example
datasets, README workflow guidance, vignettes, citation metadata, and
synthetic Gazepoint-style demonstration resources. The manuscript
demonstration uses synthetic data so that the workflow can be reproduced
publicly without redistributing private participant exports. Real
empirical studies using private Gazepoint exports should report
recruitment, ethics approval, calibration procedure, experimental
design, preprocessing decisions, exclusion rules, and confirmatory model
specifications.

\section{Limitations}\label{limitations}

The first limitation is scope. \texttt{gp3tools} is designed for
Gazepoint GP3 / Gazepoint Analysis CSV exports and should not be
presented as a universal solution for all eye-tracking platforms, export
formats, or experimental paradigms. Other devices and software
ecosystems may require different import layers, timing assumptions,
coordinate systems, calibration diagnostics, or event structures.

The second limitation concerns synthetic demonstration data. Synthetic
data make the workflow public and reproducible, but they cannot
represent every real-world artifact pattern. Real Gazepoint projects may
include calibration problems, participant movement, unusual timing
structures, inconsistent AOI labels, missing files, duplicate files,
different screen setups, stimulus-specific brightness differences, task
interruptions, or nonstandard summary exports. The demonstration is
therefore a reproducibility case study, not a comprehensive benchmark
for all possible Gazepoint data problems.

The third limitation concerns statistical interpretation.
\texttt{gp3tools} prepares model-ready data and provides representative
workflows for LMMs, GLMMs, GAMMs, time-course models, AOI analyses,
pupil-window analyses, transition summaries, and sensitivity checks.
These functions do not determine the correct scientific model
automatically. Researchers remain responsible for theoretical
justification of AOIs and windows, model specification, assumption
checking, multiplicity decisions, and the distinction between
confirmatory and exploratory analyses.

Finally, some advanced workflows depend on optional packages or
study-specific information. External-package checks require compatible
installations and input schemas; stimulus-luminance auditing requires
stimulus image files; gaze recalibration requires known target
coordinates; and transition-count sensitivity requires appropriate count
data and modelling dependencies. The package reports controlled statuses
when optional components are unavailable, but users must decide which
branches are appropriate for their study.

\section{Discussion}\label{discussion}

\texttt{gp3tools} addresses a practical reproducibility problem in
Gazepoint-based eye-tracking research. Gazepoint exports provide rich
CSV data, but researchers still need a transparent path from exported
files to analysis-ready objects, diagnostic summaries, model-ready
tables, figures, and reporting outputs. By making file validation,
sampling diagnostics, tracking-quality summaries, AOI checks, pupil
preprocessing, model preparation, and reporting readiness explicit, the
package reduces reliance on undocumented manual steps.

The contribution is workflow integration. \texttt{gp3tools} does not
introduce a new statistical model and does not claim to establish
universal preprocessing standards for eye-tracking research. Instead, it
provides Gazepoint-specific infrastructure that can be combined with
existing R modelling packages and related eye-tracking tools. This
design makes it useful for applied researchers who need to process
private Gazepoint exports while sharing reproducible analysis code,
synthetic demonstration data, and transparent diagnostic outputs.

The synthetic case study shows how a workflow can be reported without
exposing private participant data. This is important because raw
eye-tracking exports may contain sensitive participant-level quality
patterns, proprietary stimulus information, or study-specific metadata.
Public synthetic workflows allow package behaviour to be inspected while
leaving substantive empirical inference to study-specific datasets.

Future development should focus on documentation, external validation,
user-facing examples, and interoperability. Additional testing on
independent Gazepoint datasets from different laboratories, tasks,
stimulus types, screen configurations, and participant populations would
strengthen confidence in edge-case handling. Further work could also
expand export adapters, improve graphical summaries, and add examples
for common applied research designs.

\section{Supplementary material}\label{supplementary-material}

Supplementary Material S1 contains the full synthetic workflow, extended
diagnostics, code, plots, model outputs, function inventory, and
reproducibility information. The supplementary source file is
\texttt{paper-showcase/supplement/gp3tools\_full\_showcase\_supplement.Rmd}.

\section{Conclusion}\label{conclusion}

\texttt{gp3tools} provides a reproducible R workflow for Gazepoint GP3 /
Gazepoint Analysis eye-tracking exports. It bridges the gap between
vendor-specific CSV files and transparent statistical analysis by
supporting import, validation, preprocessing, AOI summaries, pupil
workflows, fixation and transition summaries, model-ready data
preparation, diagnostic plotting, reporting checklists, and archived
reproducibility materials. The package complements existing eye-tracking
and pupillometry tools by focusing on Gazepoint-specific export
structures and end-to-end workflow transparency. The synthetic
demonstration validates package behaviour and workflow reproducibility,
not substantive psychological effects. Used carefully, \texttt{gp3tools}
can make Gazepoint-based eye-tracking analyses more auditable,
reproducible, and publication-ready.

\section*{References}\label{references}
\addcontentsline{toc}{section}{References}

\protect\phantomsection\label{refs}
\begin{CSLReferences}{1}{1}
\bibitem[\citeproctext]{ref-Balaskas2026gp3tools}
Balaskas, Stefanos. 2026. \emph{Gp3tools: Import, Inspect, Analyse, and
Report Gazepoint GP3 Exports}. Zenodo.
\url{https://doi.org/10.5281/zenodo.20731340}.

\bibitem[\citeproctext]{ref-Bates2015}
Bates, Douglas, Martin Maechler, Benjamin M. Bolker, and Steve Walker.
2015. {``Fitting Linear Mixed-Effects Models Using Lme4.''}
\emph{Journal of Statistical Software} 67 (1): 1--48.
\url{https://doi.org/10.18637/jss.v067.i01}.

\bibitem[\citeproctext]{ref-Brooks2017}
Brooks, Mollie E., Kasper Kristensen, Koen J. van Benthem, et al. 2017.
{``glmmTMB Balances Speed and Flexibility Among Packages for
Zero-Inflated Generalized Linear Mixed Modeling.''} \emph{The R Journal}
9 (2): 378--400. \url{https://doi.org/10.32614/RJ-2017-066}.

\bibitem[\citeproctext]{ref-DinkFerguson2015}
Dink, Jacob W., and Brock Ferguson. 2015. \emph{eyetrackingR: An r
Library for Eye-Tracking Data Analysis}.
\url{https://www.eyetracking-r.com/}.

\bibitem[\citeproctext]{ref-Forbes2020}
Forbes, Sam H. 2020. {``PupillometryR: An r Package for Preparing and
Analysing Pupillometry Data.''} \emph{Journal of Open Source Software} 5
(50): 2285. \url{https://doi.org/10.21105/joss.02285}.

\bibitem[\citeproctext]{ref-GazepointGP3}
Gazepoint Research Inc. 2026a. \emph{Gazepoint GP3 Eye Tracker
Specifications}. \url{https://www.gazept.com/}.

\bibitem[\citeproctext]{ref-GazepointAPI}
Gazepoint Research Inc. 2026b. \emph{Open Gaze API Documentation}.
\url{https://www.gazept.com/}.

\bibitem[\citeproctext]{ref-Katz2021}
Katz, Daniel S., Morane Gruenpeter, Tom Honeyman, et al. 2021. {``A
Fresh Look at FAIR for Research Software.''} \emph{Patterns} 2 (10):
100222. \url{https://doi.org/10.1016/j.patter.2021.100222}.

\bibitem[\citeproctext]{ref-KretSjakShie2019}
Kret, Mariska E., and Elio E. Sjak-Shie. 2019. {``Preprocessing Pupil
Size Data: Guidelines and Code.''} \emph{Behavior Research Methods} 51:
1336--42. \url{https://doi.org/10.3758/s13428-018-1075-y}.

\bibitem[\citeproctext]{ref-Mathot2018Pupillometry}
Mathot, Sebastiaan. 2018. {``Pupillometry: Psychology, Physiology, and
Function.''} \emph{Journal of Cognition} 1 (1): 16.
\url{https://doi.org/10.5334/joc.18}.

\bibitem[\citeproctext]{ref-Mathot2018Baseline}
Mathot, Sebastiaan, Jasper Fabius, Ellen Van Heusden, and Stefan Van der
Stigchel. 2018. {``Safe and Sensible Preprocessing and Baseline
Correction of Pupil-Size Data.''} \emph{Behavior Research Methods} 50:
94--106. \url{https://doi.org/10.3758/s13428-017-1007-2}.

\bibitem[\citeproctext]{ref-Orquin2016}
Orquin, Jacob Lund, Nathaniel J. S. Ashby, and Alasdair D. F. Clarke.
2016. {``Areas of Interest as a Signal Detection Problem in Behavioral
Eye-Tracking Research.''} \emph{Journal of Behavioral Decision Making}
29 (2--3): 103--15. \url{https://doi.org/10.1002/bdm.1867}.

\bibitem[\citeproctext]{ref-RCoreTeam2026}
R Core Team. 2026. \emph{R: A Language and Environment for Statistical
Computing}. R Foundation for Statistical Computing.
\url{https://www.R-project.org/}.

\bibitem[\citeproctext]{ref-Smith2016}
Smith, Arfon M., Daniel S. Katz, Kyle E. Niemeyer, and FORCE11 Software
Citation Working Group. 2016. {``Software Citation Principles.''}
\emph{PeerJ Computer Science} 2: e86.
\url{https://doi.org/10.7717/peerj-cs.86}.

\bibitem[\citeproctext]{ref-Wickham2011testthat}
Wickham, Hadley. 2011. {``Testthat: Get Started with Testing.''}
\emph{The R Journal} 3: 5--10.
\url{https://journal.r-project.org/articles/RJ-2011-002/}.

\bibitem[\citeproctext]{ref-Wickham2016ggplot2}
Wickham, Hadley. 2016. \emph{Ggplot2: Elegant Graphics for Data
Analysis}. Springer. \url{https://ggplot2.tidyverse.org}.

\bibitem[\citeproctext]{ref-Wilson2014}
Wilson, Greg, D. A. Aruliah, C. Titus Brown, et al. 2014. {``Best
Practices for Scientific Computing.''} \emph{PLOS Biology} 12 (1):
e1001745. \url{https://doi.org/10.1371/journal.pbio.1001745}.

\bibitem[\citeproctext]{ref-Wood2017}
Wood, Simon N. 2017. \emph{Generalized Additive Models: An Introduction
with r}. 2nd ed. Chapman; Hall/CRC.

\bibitem[\citeproctext]{ref-Xie2015knitr}
Xie, Yihui. 2015. \emph{Dynamic Documents with r and Knitr}. 2nd ed.
Chapman; Hall/CRC.

\bibitem[\citeproctext]{ref-XieAllaireGrolemund2018}
Xie, Yihui, J. J. Allaire, and Garrett Grolemund. 2018. \emph{R
Markdown: The Definitive Guide}. Chapman; Hall/CRC.
\url{https://bookdown.org/yihui/rmarkdown/}.

\end{CSLReferences}

\end{document}
